\documentclass[12pt,a4paper]{article}
\usepackage[utf8]{inputenc}
\usepackage[T1]{fontenc}
\usepackage{lmodern}
\usepackage{amsmath,amsfonts,amssymb}
\usepackage{graphicx}
\usepackage[colorlinks=true,linkcolor=blue,citecolor=blue,urlcolor=blue]{hyperref}
\usepackage{booktabs}
\usepackage{multirow}
\usepackage{geometry}
\usepackage{setspace}
\usepackage[numbers,sort&compress]{natbib}
\usepackage{float}
\usepackage{caption}
\usepackage{subcaption}
\usepackage{authblk}
\usepackage{parskip}
\geometry{a4paper, top=1in, bottom=1in, left=1.2in, right=1.2in}
\setstretch{1.2}

\title{\textbf{Predicting Citation Links in Library and Information Science:\\
A Temporally-Aware Machine Learning Framework\\
with Network, Bibliometric, and Semantic Features}}

\author[1]{Moses Boudourides}
\author[2]{Giannis Tsakonas}
\affil[1]{Department of Mathematics, University of Patras, Greece\\
  \texttt{m.boudourides@upatras.gr}}
\affil[2]{Library \& Information Center, University of Patras, Greece\\
  \texttt{gtsakonas@upatras.gr}}
\date{\today}

\begin{document}
\maketitle
\thispagestyle{empty}

\begin{abstract}
\noindent
The prediction of citation links between scientific papers is a fundamental problem
in computational bibliometrics and the sociology of science. While preferential
attachment and cumulative advantage explain the macroscopic topology of citation
networks, predicting microscopic, pair-level citations requires integrating
structural, temporal, semantic, and bibliometric signals in a causally rigorous
framework. In this study, we formulate citation link prediction as a binary
classification task on two large-scale, independent datasets in Library and
Information Science (LIS): 259,220 articles from Dimensions.ai and 168,901 articles
from OpenAlex, both spanning 1975--2024.

To prevent temporal data leakage---a pervasive flaw in existing citation prediction
literature---we enforce strict temporal causality throughout: all features for a
potential citation from paper $A$ (published in year $t_A$) to paper $B$ are
computed exclusively from data available prior to year $t_A$. We employ a rigorous
hard-negative sampling strategy, pairing each positive citation with a non-citation
drawn from the same temporal window ($\pm 3$ years) and fine-grained disciplinary
category, forcing models to learn genuine citation dynamics rather than trivial
temporal or topical differences.

We engineer a rich feature set spanning five categories: prestige (year-capped
citation count and PageRank), temporal dynamics (citation age and in-degree
velocity), bibliographic coupling (common references, Jaccard similarity, co-citation
overlap), semantic proximity (TF-IDF cosine similarity), and bibliometric metadata
(journal prestige, author productivity). Using Gradient Boosting, we achieve
ROC-AUC of 0.8949 $\pm$ 0.0006 (Dimensions) and 0.9762 $\pm$ 0.0004 (OpenAlex) under 5-fold stratified cross-validation, with Brier scores of 0.1346 and 0.0871 respectively, confirming well-calibrated probability estimates. McNemar's test confirms statistical superiority of Gradient Boosting over Random Forest ($p < 10^{-13}$ for both datasets). Rolling temporal evaluation across four independent test windows confirms stable performance over two decades.

A dedicated semantic ablation experiment reveals that semantic proximity is the
dominant single-feature predictor in OpenAlex (AUC drop of 0.0269 when removed), while temporal gap and co-citation overlap dominate in the purely structural Dimensions setting. We discuss the epistemological implications of these findings, arguing that the high predictability of citations reflects the bounded, paradigmatic nature of knowledge diffusion within established disciplinary communities rather than a trivial methodological artifact.
\end{abstract}

\newpage
\tableofcontents
\newpage

%% ============================================================
\section{Introduction}
\label{sec:intro}
%% ============================================================

The act of citation is the fundamental mechanism through which scientific knowledge
is validated, contested, and integrated into the scholarly record
\citep{merton1968matthew, price1965networks}. Citations form a vast, directed
network that encodes both the intellectual lineage of individual papers and the
social structure of scientific communities \citep{fortunato2018science}. From a
macroscopic perspective, the evolution of citation networks is well-described by
generative models such as preferential attachment \citep{barabasi1999emergence,
de2000preferential}: highly cited papers accumulate new citations at a rate
proportional to their existing prestige, generating the characteristic power-law
distribution of citation counts observed across all scientific fields. However,
while these models successfully reproduce aggregate network topology, they are
insufficient for microscopic prediction: given a newly published paper $A$ and a
previously published paper $B$, what is the probability that $A$ will cite $B$?

This pair-level prediction task is central to both the sociology of science and the
design of modern scientific recommender systems \citep{shibata2012link,
bhagavatula2018content}. The challenge lies in the fact that citations are not
purely random nor purely prestige-driven; they are deliberate acts of intellectual
positioning constrained by temporal relevance, topical proximity, and social
structure \citep{kuhn1962structure, bornmann2008what}. A researcher choosing which
papers to cite does not merely select the most highly cited works in a field; they
select works that are relevant to their specific research question, known to them
through their social and institutional networks, and established as credible within
their disciplinary community \citep{crane1972invisible}.

In recent years, machine learning approaches have increasingly been applied to
citation link prediction. However, as noted by several critics, many such models
suffer from subtle methodological flaws that artificially inflate performance
\citep{wang2013quantifying}. The two most pervasive issues are \textit{temporal
data leakage}---using future information, such as the total static citation count
of paper $B$, to predict a citation from paper $A$ published years earlier---and
\textit{trivial negative sampling}---pairing a recent paper with a random, entirely
unrelated paper from decades ago or from a different field, making the classification
task artificially easy. These flaws are not merely technical inconveniences; they
undermine the validity of any theoretical claims derived from inflated performance
metrics.

In this study, we address these methodological challenges directly. We present a
strictly temporally-aware machine learning pipeline for predicting citation links
within the domain of Library and Information Science (LIS), evaluated on two
comprehensive and independent datasets derived from Dimensions.ai and OpenAlex.
LIS is an ideal testbed for this investigation: it is a mature, interdisciplinary
field that has itself contributed substantially to bibliometric methodology, yet
its own citation dynamics have received comparatively less attention than the
physical or biomedical sciences. Its breadth---spanning information retrieval,
knowledge organization, scholarly communication, and digital libraries---makes it
a rich and heterogeneous corpus for studying citation behavior across diverse
research traditions.

Our contributions are fourfold. First, we introduce a rigorous, temporally-aware
evaluation methodology that substantially mitigates data leakage by capping all
network and prestige features strictly prior to the citing paper's publication year.
Second, we enrich the standard bibliometric feature set with genuine graph-theoretic
features (year-capped PageRank, in-degree velocity) and bibliometric metadata
(journal prestige, author productivity), demonstrating their incremental predictive
value via leave-one-out ablation. Third, we conduct a dedicated semantic ablation
experiment to quantify the independent contribution of semantic proximity and to
directly address the concern that high performance in semantically-constrained
settings may reflect topical redundancy rather than genuine citation dynamics.
Fourth, we discuss the epistemological implications of high citation predictability
for the sociology of science, connecting our empirical findings to broader debates
about the nature of knowledge diffusion, disciplinary closure, and the structural
determination of scientific attention.

The remainder of this paper is organized as follows. Section~\ref{sec:litreview}
reviews the relevant literature. Section~\ref{sec:data} describes our data sources,
feature engineering, and experimental setup. Sections~\ref{sec:dim}
and~\ref{sec:oa} present results for the Dimensions and OpenAlex datasets.
Section~\ref{sec:comp} provides a comparative analysis. Section~\ref{sec:discussion}
discusses theoretical and methodological implications. Section~\ref{sec:conclusion}
concludes.

%% ============================================================
\section{Literature Review}
\label{sec:litreview}
%% ============================================================

\subsection{Classical Theories of Citation Dynamics}

The study of citation patterns is a cornerstone of bibliometrics and the science
of science. The earliest and most enduring model of citation dynamics is Merton's
concept of the Matthew Effect \citep{merton1968matthew}: the tendency for already
recognized scientists and papers to receive disproportionate future recognition,
generating a ``rich-get-richer'' dynamic. De Solla Price formalized this as
cumulative advantage in citation networks \citep{de1976general}, and Barab\'asi
and Albert later generalized it as preferential attachment---a universal mechanism
for the emergence of scale-free networks \citep{barabasi1999emergence}. These
models provide a compelling account of the aggregate statistical properties of
citation networks but are insufficient for pair-level prediction.

Subsequent research has demonstrated that the strength of preferential attachment
varies significantly across disciplines and over time \citep{wang2008measuring}.
Bianco and Gabrielli introduced the concept of paper ``fitness,'' suggesting that
intrinsic quality or relevance modulates the rich-get-richer effect
\citep{bianco2008fitness}. Wang, Song, and Barab\'asi demonstrated that the
long-term scientific impact of a paper is governed by a combination of initial
fitness and cumulative advantage, challenging purely prestige-based models
\citep{wang2013quantifying}. These theoretical developments motivate our
multi-feature approach: citation behavior is not a simple function of prestige
but emerges from the complex interaction of prestige, relevance, and social structure.

\subsection{Social and Semantic Proximity in Citation}

Beyond prestige, a growing body of research has demonstrated that social proximity
and semantic proximity are powerful predictors of citation behavior, often rivaling
or exceeding the effect of prestige for the majority of papers. Newman's foundational
work on co-authorship networks established that scientific collaboration is highly
clustered and that co-authorship strongly predicts future collaboration and citation
\citep{newman2004coauthorship}. Kumar extended this analysis to show that citation
probability decays rapidly with social distance in co-authorship networks
\citep{kumar2015co}. Martin et al.\ demonstrated that the social structure of
scientific communities---operationalized through co-authorship, institutional
affiliation, and conference participation---is a major determinant of citation
patterns \citep{martin2013coauthorship}.

The landmark ``Citation proximus'' study by Kozlowski et al.\ \citep{kozlowski2025citation}
provided rigorous statistical evidence for the joint role of social and semantic
proximity using Generalized Linear Mixed Models, demonstrating that these factors
together explain a substantial fraction of the variance in citation behavior across
multiple scientific fields. This study opened the door for more computationally
advanced, predictive inquiry: the relationships between network structure, semantic
content, and prestige are unlikely to be simple and additive; they involve complex,
non-linear interactions that traditional statistical models may fail to capture fully.

Semantic proximity has been operationalized through a range of methods, from
keyword overlap and topic models \citep{blei2003latent} to deep neural embeddings
\citep{devlin2019bert, reimers2019sentence}. While contextual embeddings from
models such as BERT and SBERT offer greater expressiveness, they introduce a
potential leakage concern in citation prediction: pre-trained language models may
have been exposed to the target corpus during their training phase, potentially
inflating similarity scores for papers that appear in their training data. We
address this concern by using TF-IDF cosine similarity, computed solely from the
corpus at hand, which is free from this concern while remaining a well-validated
measure of topical proximity.

\subsection{Temporal Dynamics and Citation Obsolescence}

A dimension often underappreciated in citation prediction is the temporal decay
of citation probability. Price's early work on citation half-lives established
that the probability of a paper being cited declines with age, reflecting both
the natural obsolescence of older work and the tendency for researchers to
prioritize recent literature \citep{de1976general}. Redner's analysis of citation
distributions confirmed that temporal decay is highly non-uniform: foundational
papers accumulate citations over decades, while applied or methodological papers
may be superseded within a few years \citep{redner2005citation}. The temporal gap
between citing and cited paper is therefore a complex predictor that interacts
with prestige, field, and the nature of the contribution.

Temporal link prediction---predicting which edges will appear in a network at
future time steps---is a well-studied problem in network science
\citep{liben2007link}. In citation networks, temporal prediction is particularly
challenging because the network is acyclic (a paper can only cite papers published
before it) and because the candidate space grows quadratically with corpus size.
Shibata et al.\ applied temporal link prediction methods to citation networks,
finding that structural features alone can achieve reasonable predictive performance
but that temporal constraints are essential for valid evaluation \citep{shibata2012link}.
Our rolling temporal evaluation framework directly addresses this challenge by
training on papers published before a cutoff year and testing on papers published
in a subsequent window, providing a rigorous assessment of temporal generalizability.

\subsection{Bibliographic Coupling and Co-citation Analysis}

Beyond prestige and temporal dynamics, the structural overlap between reference
lists provides a powerful signal for citation prediction. Bibliographic coupling
\citep{kessler1963bibliographic}---the number of references shared by two
papers---reflects shared intellectual heritage. Co-citation analysis
\citep{small1973co}---measuring how often two papers are cited together by a
third paper---captures a complementary dimension of intellectual proximity.
Together, these reference overlap features capture the bibliographic embedding of
a paper in the existing literature, a dimension of proximity that is distinct from
both semantic content and social network position.

The Jaccard coefficient of reference sets provides a normalized measure of
bibliographic coupling that is invariant to differences in reference list length.
Common citers (the number of papers that cite both the citing and cited paper,
excluding the citing paper itself) operationalize co-citation overlap in a
temporally causal manner. We note that for recently published papers with few
citations, co-citation overlap will be sparse or zero; this is an inherent property
of the measure rather than a methodological artifact, and it is captured implicitly
by the model through its interaction with prestige and temporal gap features.

\subsection{Machine Learning for Citation Prediction}

The integration of machine learning into bibliometrics has shifted the focus from
explanatory models to predictive frameworks. Early predictive efforts employed
standard regression models to predict citation counts \citep{lokker2008prediction}.
More recent studies have leveraged ensemble methods such as Gradient Boosting
\citep{natekin2013gradient} and Random Forests \citep{breiman2001random}, which
are particularly adept at capturing the non-linear interactions between disparate
feature types. Chakraborty et al.\ proposed a stratified learning approach for
citation count prediction, demonstrating that models trained on subgroups outperform
models trained on the full dataset \citep{chakraborty2015towards}. Bhagavatula
et al.\ applied content-based citation recommendation using neural embeddings,
demonstrating the value of semantic features for citation prediction
\citep{bhagavatula2018content}. Alohali et al.\ provided a recent systematic
review of machine learning approaches to citation analysis \citep{alohali2022machine}.

The interpretability of machine learning models is a growing concern in
bibliometrics. SHAP (SHapley Additive exPlanations) \citep{lundberg2017unified}
provides a theoretically grounded framework for attributing model predictions to
individual features, enabling researchers to identify which factors drive citation
decisions in a given corpus. We employ SHAP analysis throughout this study to
bridge the gap between predictive performance and theoretical interpretation.

\subsection{Bibliographic Data Infrastructure}

Dimensions.ai \citep{herzog2020dimensions} is a comprehensive, linked research
information system integrating publications, citations, grants, patents, and policy
documents. Its coverage of the scientific literature is among the broadest
available, and it provides structured access to metadata including author
affiliations, reference lists, and Field of Research (FoR) classifications based
on the ANZSRC 2020 taxonomy.

OpenAlex \citep{priem2022openalex} is a fully open, freely accessible index of
scholarly works developed as a successor to Microsoft Academic Graph. OpenAlex
is particularly notable for its open data philosophy and its rich topic taxonomy,
which enables fine-grained subject classification. The use of two independent data
sources in this study is not merely a robustness check; it is a substantive
methodological contribution, as findings consistent across Dimensions and OpenAlex
are more likely to reflect genuine properties of LIS citation dynamics rather than
artifacts of a specific data infrastructure.

%% ============================================================
\section{Data and Methodology}
\label{sec:data}
%% ============================================================

\subsection{Corpus Construction}

\subsubsection{Dimensions.ai Dataset}

The Dimensions corpus was constructed by querying the Dimensions.ai API for all
journal articles in Library and Information Science published between 1975 and
2024. The retrieval strategy combined two complementary approaches to maximize
coverage while maintaining topical relevance.

The primary query filtered by Field of Research (FoR) classification code
\texttt{4610} (Library and Information Studies) under the ANZSRC 2020 taxonomy,
executed year by year from 1975 to 2024:

\begin{verbatim}
search publications
where category_for.id = "4610"
and year = {Y}
and type = "article"
return publications[id+year+title+authors+journal+
  abstract+times_cited+reference_ids+category_for+doi]
\end{verbatim}

To supplement this with articles that may have been misclassified or assigned to
adjacent FoR codes, four additional keyword queries were executed against titles
and abstracts for each year:

\begin{verbatim}
search publications in title_abstract_only
for "{keyword}"
where year = {Y} and type = "article"
\end{verbatim}

\noindent where \texttt{keyword} $\in$ \{\textit{knowledge organization},
\textit{digital libraries}, \textit{information literacy},
\textit{academic libraries}\}. These keywords were selected to capture the major
LIS subdisciplines most likely to appear in adjacent FoR categories (education,
computer science). Results were deduplicated by Dimensions article ID. The
retrieved fields included: unique identifier, publication year, full title, author
list with Dimensions author IDs, journal name, abstract text, citation count,
reference ID list, FoR classification codes, and DOI. The final Dimensions corpus
comprises \textbf{259,220 articles}, of which 96.3\% have abstracts available.

\subsubsection{OpenAlex Dataset}

The OpenAlex corpus was constructed using the OpenAlex REST API, filtering by five
specific \texttt{primary\_topic.id} values corresponding to core LIS and
scientometrics research areas. The legacy OpenAlex concept \texttt{C136764020}
(``Library and Information Science'') was deliberately avoided, as it is assigned
by an automated ML tagger and returns approximately 4.2 million records---far too
broad for a focused citation dynamics study. The five topic IDs used were:

\begin{center}
\begin{tabular}{ll}
\toprule
\textbf{Topic ID} & \textbf{Description} \\
\midrule
\texttt{T10712} & Library Science and Information Literacy \\
\texttt{T14330} & Library Science and Information Systems \\
\texttt{T13166} & Information Science and Libraries \\
\texttt{T13673} & Library Science and Information \\
\texttt{T10102} & Scientometrics and Bibliometrics Research \\
\bottomrule
\end{tabular}
\end{center}

The filter applied OR logic across all five topic IDs, restricted to article type
and publication years 1975--2024. The API was queried using the polite pool with
an institutional email address. After deduplication by OpenAlex work ID, the final
OpenAlex corpus comprises \textbf{168,901 articles}, of which 64.3\% have abstracts
available.

\subsubsection{Corpus Characteristics}

\begin{table}[H]
\centering
\caption{Corpus Characteristics}
\label{tab:corpus}
\begin{tabular}{lcc}
\toprule
\textbf{Characteristic} & \textbf{Dimensions} & \textbf{OpenAlex} \\
\midrule
Total articles & 259,220 & 168,901 \\
Year range & 1975--2024 & 1975--2024 \\
Articles with abstracts & 249,751 (96.3\%) & 108,611 (64.3\%) \\
Mean citations per article & 20.98 & 6.26 \\
Median citations per article & 4 & 0 \\
Within-corpus citation pairs & 331,047 & 239,349 \\
\bottomrule
\end{tabular}
\end{table}

\subsection{Citation Pair Construction and Hard Negative Sampling}

The fundamental task is binary classification: given a pair of papers $(A, B)$
where $t_B \le t_A$, predict whether $A$ cites $B$. We extracted all observed
within-corpus citations where both the citing paper $A$ and the cited paper $B$
exist in our dataset, yielding 331,047 positive pairs for Dimensions and 239,349
for OpenAlex.

For each positive pair $(A, B)$, we generated a hard negative pair $(A, B')$ by
sampling a paper $B'$ that satisfies all of the following constraints:
\begin{enumerate}
  \item $|t_{B'} - t_B| \le 3$ (temporal plausibility);
  \item $B'$ shares at least one Field of Research code with $B$ (Dimensions) or
    belongs to the same primary topic cluster (OpenAlex) (topical comparability);
  \item $A$ does not actually cite $B'$ (verified against the full reference list);
  \item $B' \ne B$.
\end{enumerate}

This hard negative sampling strategy is substantially more demanding than uniform
random sampling. By constraining negatives to the same temporal neighborhood and
topical domain as the positive cited paper, we ensure that the classification task
cannot be solved by trivially exploiting temporal or topical differences. The
resulting datasets contain \textbf{662,094 balanced pairs} for Dimensions and
\textbf{478,698 balanced pairs} for OpenAlex.

\subsection{Temporally-Aware Feature Engineering}

We engineer features in five categories. All features that depend on the state of
the literature are computed using a strictly causal rolling window: for any citing
paper published in year $t$, every feature is derived solely from papers published
up to year $t-1$.

\subsubsection{Prestige Features}

\textbf{Prestige of Cited Paper ($P_{cited}$):} The log-transformed cumulative
citation count received by paper $B$ from papers published up to year $t_A - 1$:
\begin{equation}
P_{cited}(B, t_A) = \log\!\left(1 + \sum_{C\,:\,t_C \le t_A - 1} \mathbf{1}[C \text{ cites } B]\right)
\end{equation}
This is the primary operationalization of the Matthew Effect: papers that have
already attracted many citations are more likely to attract future ones.

\textbf{Activity of Citing Paper ($A_{citing}$):} The log-transformed cumulative
citation count of paper $A$ up to year $t_A - 1$ (OpenAlex only). This captures
the prestige of the citing paper itself.

\subsubsection{Graph-Theoretic Network Features}

\textbf{Year-Capped PageRank ($PR_{cited}$):} The PageRank centrality of paper $B$
in the directed citation graph $G_{t_A - 1}$, computed using the power-iteration
algorithm with damping factor $\alpha = 0.85$:
\begin{equation}
PR(B) = \frac{1-\alpha}{|V|} + \alpha \sum_{C \to B} \frac{PR(C)}{d^+(C)}
\end{equation}
where $d^+(C)$ is the out-degree of paper $C$ in $G_{t_A - 1}$. Unlike raw
citation count, PageRank accounts for the quality of citing papers.

\textbf{In-degree Velocity ($V_{cited}$):} The proportion of paper $B$'s total
citations (up to $t_A - 1$) received in the three years immediately preceding $t_A$:
\begin{equation}
V_{cited}(B, t_A) = \frac{\sum_{C\,:\,t_A - 3 \le t_C \le t_A - 1} \mathbf{1}[C \text{ cites } B]}{\max\!\left(1,\, P_{cited}(B, t_A)\right)}
\end{equation}
This captures the momentum of a paper's citation trajectory, distinguishing papers
experiencing a recent surge from those with steady long-term impact.

\subsubsection{Bibliographic Coupling and Co-citation Features}

\textbf{Common References ($CR$):} The number of papers cited by both $A$ and $B$:
$CR(A, B) = |R_A \cap R_B|$.

\textbf{Jaccard Reference Similarity ($J_{refs}$):} The Jaccard coefficient of
the reference sets: $J_{refs}(A, B) = |R_A \cap R_B| / |R_A \cup R_B|$.

\textbf{Co-citation Overlap ($CC$):} The number of papers (other than $A$)
published up to $t_A - 1$ that cite both $A$ and $B$:
\begin{equation}
CC(A, B) = |\{C \ne A : C \text{ cites } A \text{ and } C \text{ cites } B,\; t_C \le t_A - 1\}|
\end{equation}

\subsubsection{Temporal Feature}

\textbf{Temporal Gap ($\Delta t$):} $\Delta t(A, B) = t_A - t_B$. This reflects
the natural decay of citation probability with age and the tendency for recent
papers to cite other recent papers.

\subsubsection{Semantic Proximity}

\textbf{Semantic Similarity ($S_{sem}$):} The TF-IDF cosine similarity between
the concatenated title and abstract of papers $A$ and $B$. We fit a TF-IDF
vectorizer with up to 50,000 features and bigrams on the full corpus, then compute
pairwise cosine similarities for all citation pairs. We explicitly chose TF-IDF
over pre-trained neural embeddings (e.g., SBERT) to guarantee no temporal data
leakage: pre-trained language models may have observed the target corpus during
their training phase, potentially inflating similarity scores. TF-IDF, computed
solely from the corpus at hand, is free from this concern.

\subsubsection{Bibliometric Metadata Features (Dimensions Only)}

\textbf{Journal Prestige ($JP$):} The total citation count accumulated by all
papers published in the same journal as paper $B$, up to year $t_A - 1$.

\textbf{Author Productivity ($AP$):} The total number of papers published by the
most prolific author of paper $A$ up to year $t_A - 1$.

\subsection{Feature Summary}

\begin{table}[H]
\centering
\caption{Complete Feature Set}
\label{tab:features}
\begin{tabular}{llcc}
\toprule
\textbf{Category} & \textbf{Feature} & \textbf{Dimensions} & \textbf{OpenAlex} \\
\midrule
Prestige & $P_{cited}$ (log citation count) & \checkmark & \checkmark \\
Prestige & $A_{citing}$ (log citing activity) & -- & \checkmark \\
Graph & $PR_{cited}$ (year-capped PageRank) & \checkmark & \checkmark \\
Graph & $V_{cited}$ (in-degree velocity) & \checkmark & \checkmark \\
Bibliographic & $CR$ (common references) & \checkmark & \checkmark \\
Bibliographic & $J_{refs}$ (Jaccard reference sim.) & \checkmark & \checkmark \\
Bibliographic & $CC$ (co-citation overlap) & \checkmark & \checkmark \\
Temporal & $\Delta t$ (temporal gap) & \checkmark & \checkmark \\
Semantic & $S_{sem}$ (TF-IDF cosine sim.) & \checkmark & \checkmark \\
Metadata & $JP$ (journal prestige) & \checkmark & -- \\
Metadata & $AP$ (author productivity) & \checkmark & -- \\
\midrule
\textbf{Total} & & \textbf{9} & \textbf{9} \\
\bottomrule
\end{tabular}
\end{table}

\subsection{Machine Learning Models}

We evaluate four models spanning the linear-to-non-linear spectrum:

\textbf{Logistic Regression (LR):} A linear baseline that models the log-odds of
citation as a linear function of features. It provides a transparent benchmark
against which to evaluate the added value of non-linear models.

\textbf{Linear SVM:} A linear support vector machine that finds the maximum-margin
hyperplane separating cited from non-cited pairs, optimizing hinge loss.

\textbf{Random Forest (RF):} An ensemble of 200 decision trees trained on bootstrap
samples \citep{breiman2001random}. At each node, a random subset of $\sqrt{d}$
features is considered for splitting.

\textbf{Gradient Boosting (GB):} A sequential ensemble method that builds trees
to correct the errors of previous trees \citep{natekin2013gradient}. We use 200
trees with learning rate 0.1 and maximum depth 5.

\subsection{Evaluation Framework}

All models are evaluated using \textbf{5-fold stratified cross-validation},
preserving the class distribution in each fold. We report mean and standard
deviation across folds for six metrics: ROC-AUC, PR-AUC, F1-score, Accuracy, MCC,
and Brier score. Additionally, we conduct:
\begin{itemize}
  \item \textbf{Rolling temporal evaluation} across four independent test windows
    (train on papers before window start, test on papers within window);
  \item \textbf{Temporal hold-out} (train $\le 2015$, test 2016--2020);
  \item \textbf{McNemar's test} comparing RF and GB predictions on the full dataset;
  \item \textbf{Calibration analysis} using Brier score and reliability diagram;
  \item \textbf{SHAP analysis} and leave-one-out ablation for the best model;
  \item \textbf{Semantic ablation experiment} (OpenAlex only): full CV with and
    without semantic similarity for all four models.
\end{itemize}

%% ============================================================
\section{Results: Dimensions.ai Analysis}
\label{sec:dim}
%% ============================================================

\subsection{Corpus Overview}

Figure~\ref{fig:dim_pubs} shows the temporal distribution of LIS publications in
the Dimensions dataset. The corpus exhibits steady growth from the late 1980s
through the 2010s, reflecting both the expansion of LIS as a field and the
increasing coverage of Dimensions over time. The distribution of citation counts
is highly skewed, with a small number of highly cited papers and a large majority
with few citations---a characteristic signature of preferential attachment and the
Matthew Effect.

\begin{figure}[H]
\centering
\includegraphics[width=0.78\textwidth]{../figures/dimensions/fig1_publications_per_year.png}
\caption{Number of LIS publications per year in the Dimensions dataset (1975--2024).}
\label{fig:dim_pubs}
\end{figure}

\begin{figure}[H]
\centering
\includegraphics[width=0.78\textwidth]{../figures/dimensions/fig2_mean_citations_per_year.png}
\caption{Mean citations per publication year in the Dimensions dataset.}
\label{fig:dim_cites}
\end{figure}

Figure~\ref{fig:dim_cites} shows the mean citation count per publication year.
Papers published in the 1990s and early 2000s have accumulated the highest mean
citation counts, reflecting the combined effect of their age (more time to
accumulate citations) and their foundational role in establishing the modern LIS
research agenda. The decline in mean citations for more recent papers is expected,
as these papers have had less time to accumulate citations.

\subsection{Model Performance}


Table~\ref{tab:dim_cv} presents the 5-fold stratified cross-validation results
for all four models on the Dimensions dataset (662,094 pairs, 9 features).
Gradient Boosting achieves the highest performance across all metrics, with
ROC-AUC of $0.8949 \pm 0.0006$,
F1 of $0.8013 \pm 0.0009$,
and MCC of $0.6175 \pm 0.0016$.
The Brier score of 0.1346 confirms well-calibrated probability
estimates. McNemar's test confirms that Gradient Boosting is statistically superior
to Random Forest ($p = 7.20e-26$).


\begin{table}[H]
\centering
\caption{5-Fold CV Results: Dimensions Dataset (Mean $\pm$ Std)}
\label{tab:dim_cv}
\resizebox{\textwidth}{!}{
\begin{tabular}{lcccccc}
\toprule
\textbf{Model} & \textbf{ROC-AUC} & \textbf{PR-AUC} & \textbf{F1} & \textbf{Accuracy} & \textbf{MCC} & \textbf{Brier} \\
\midrule

Logistic Regression & $0.8719\pm0.0006$ & $0.8893\pm0.0008$ & $0.7757\pm0.0012$ & $0.7888\pm0.0009$ & $0.5816\pm0.0016$ & $0.1461\pm0.0005$ \\
Linear SVM & $0.8677\pm0.0005$ & $0.8816\pm0.0006$ & $0.7701\pm0.0008$ & $0.7841\pm0.0005$ & $0.5726\pm0.0010$ & $0.1518\pm0.0004$ \\
Random Forest & $0.8741\pm0.0010$ & $0.8858\pm0.0009$ & $0.7882\pm0.0012$ & $0.7918\pm0.0009$ & $0.5840\pm0.0018$ & $0.1477\pm0.0007$ \\
\textbf{Gradient Boosting} & $0.8949\pm0.0006$ & $0.9063\pm0.0007$ & $0.8013\pm0.0009$ & $0.8080\pm0.0008$ & $0.6175\pm0.0016$ & $0.1307\pm0.0004$ \\
\bottomrule
\end{tabular}}
\end{table}

\begin{figure}[H]
\centering
\includegraphics[width=0.85\textwidth]{../results/v2/figures/dim/dim_cv_v3.png}
\caption{Model performance comparison on the Dimensions dataset (5-fold CV).}
\label{fig:dim_cv}
\end{figure}

\subsection{SHAP Feature Importance}

Figure~\ref{fig:dim_shap} and Table~\ref{tab:dim_shap} present the SHAP feature
importance values for the Gradient Boosting model on the Dimensions dataset.
The dominant predictor is \texttt{prestige\_cited}
(mean $|$SHAP$|$ = 0.8105), followed by
\texttt{common\_citers} (0.5125) and
\texttt{indegree\_velocity} (0.4954).
The prestige of the cited paper is the strongest individual predictor, confirming
the Matthew Effect: papers that have already accumulated citations are more likely
to attract further citations. The co-citation overlap and in-degree velocity features
rank second and third, reflecting the importance of momentum and shared intellectual
heritage. The graph-theoretic features (PageRank, in-degree velocity) contribute
meaningfully beyond raw citation count, confirming the value of the enriched feature
set. Journal prestige and author productivity, while individually modest, each
contribute positively to the model.

\begin{table}[H]
\centering
\caption{SHAP Feature Importance: Dimensions Dataset (Gradient Boosting)}
\label{tab:dim_shap}
\begin{tabular}{lcc}
\toprule
\textbf{Feature} & \textbf{Mean $|$SHAP$|$} & \textbf{Rank} \\
\midrule

prestige\_cited & 0.8105 & 1 \\
common\_citers & 0.5125 & 2 \\
indegree\_velocity & 0.4954 & 3 \\
temporal\_gap & 0.4293 & 4 \\
jaccard\_refs & 0.3710 & 5 \\
pagerank\_cited & 0.2428 & 6 \\
journal\_prestige & 0.1793 & 7 \\
author\_productivity & 0.1430 & 8 \\
common\_refs & 0.1215 & 9 \\
\bottomrule
\end{tabular}
\end{table}

\begin{figure}[H]
\centering
\includegraphics[width=0.75\textwidth]{../results/v2/figures/dim/dim_shap_v3.png}
\caption{SHAP feature importance for Gradient Boosting on the Dimensions dataset.}
\label{fig:dim_shap}
\end{figure}

\subsection{Ablation Study}

Table~\ref{tab:dim_ablation} presents the leave-one-out ablation study results
for the Gradient Boosting model on the Dimensions dataset. The baseline AUC
(all features) is 0.8945. Removing the temporal gap causes
the largest AUC drop (0.0110), followed by
co-citation overlap (0.0098) and prestige
(0.0061). The graph-theoretic features
(PageRank, in-degree velocity) each contribute positively, validating the enriched
feature engineering approach.

\begin{table}[H]
\centering
\caption{Leave-One-Out Ablation: Dimensions Dataset (Gradient Boosting)}
\label{tab:dim_ablation}
\begin{tabular}{lcc}
\toprule
\textbf{Feature Removed} & \textbf{AUC Without} & \textbf{AUC Drop} \\
\midrule

temporal\_gap & 0.8835 & $-0.0110$ \\
common\_citers & 0.8847 & $-0.0098$ \\
prestige\_cited & 0.8884 & $-0.0061$ \\
journal\_prestige & 0.8909 & $-0.0036$ \\
author\_productivity & 0.8928 & $-0.0018$ \\
pagerank\_cited & 0.8933 & $-0.0012$ \\
indegree\_velocity & 0.8941 & $-0.0005$ \\
jaccard\_refs & 0.8942 & $-0.0003$ \\
common\_refs & 0.8945 & $-0.0000$ \\
\bottomrule
\end{tabular}
\end{table}

\begin{figure}[H]
\centering
\includegraphics[width=0.75\textwidth]{../results/v2/figures/dim/dim_ablation_v3.png}
\caption{Leave-one-out ablation study for Gradient Boosting on the Dimensions dataset.}
\label{fig:dim_ablation}
\end{figure}

\subsection{Rolling Temporal Evaluation}

Table~\ref{tab:dim_rolling} presents the rolling temporal evaluation results for
all four models across four independent test windows. Performance is stable across
all periods, with Gradient Boosting consistently achieving the highest AUC. The
modest decline in the 2018--2025 window reflects the natural difficulty of
predicting citations for very recent papers, which have had less time to accumulate
the structural signals (prestige, co-citation) that drive the model.

\begin{table}[H]
\centering
\caption{Rolling Temporal Evaluation: Dimensions Dataset (ROC-AUC)}
\label{tab:dim_rolling}
\begin{tabular}{lcccc}
\toprule
\textbf{Test Window} & \textbf{LR} & \textbf{SVM} & \textbf{RF} & \textbf{GB} \\
\midrule

2005-2010 & 0.8665 & 0.8615 & 0.8612 & 0.8732 \\
2010-2015 & 0.8754 & 0.8674 & 0.8729 & 0.8901 \\
2015-2020 & 0.8761 & 0.8684 & 0.8752 & 0.8948 \\
2018-2025 & 0.8675 & 0.8595 & 0.8636 & 0.8712 \\
\bottomrule
\end{tabular}
\end{table}

\begin{figure}[H]
\centering
\includegraphics[width=0.85\textwidth]{../results/v2/figures/dim/dim_rolling_v3.png}
\caption{Rolling temporal evaluation for all models on the Dimensions dataset.}
\label{fig:dim_rolling}
\end{figure}

\subsection{Temporal Hold-Out Evaluation}

As an additional robustness check, we trained all models on pairs with citing year
$\le 2015$ and tested on pairs with citing year 2016--2020. Gradient Boosting
achieves AUC of 0.8911 on this out-of-time test
set, compared to 0.8658 for Logistic Regression,
confirming that the model generalizes robustly across time periods and that the
performance gains of non-linear models are not an artifact of overfitting to the
training period.


%% ============================================================
\section{Results: OpenAlex Analysis}
\label{sec:oa}
%% ============================================================

\subsection{Corpus Overview}

\begin{figure}[H]
\centering
\includegraphics[width=0.78\textwidth]{../figures/oa/oa_fig1_publications_per_year.png}
\caption{Number of LIS publications per year in the OpenAlex dataset (1975--2024).}
\label{fig:oa_pubs}
\end{figure}

The OpenAlex corpus exhibits a similar growth trajectory to Dimensions, with strong
growth from the 1990s onward. The lower mean citation count (6.26 vs.\ 20.98 for
Dimensions) reflects both the broader topical scope of the OpenAlex corpus and the
lower baseline citation rates for open-access articles, which constitute a larger
fraction of the OpenAlex corpus.

\subsection{Model Performance}


Table~\ref{tab:oa_cv} presents the 5-fold cross-validation results for the
OpenAlex dataset (478,698 pairs, 9 features). Gradient Boosting achieves ROC-AUC
of $0.9762 \pm 0.0004$, F1 of
$0.9219 \pm 0.0009$, and Brier score of
0.0871. McNemar's test confirms statistical superiority of
Gradient Boosting over Random Forest ($p = 4.93e-14$).

It is important to contextualize these metrics. While the AUC values are high, they
are obtained on a balanced binary classification task where negative pairs are
constrained to the same temporal and topical neighborhood as positive pairs. The
real-world citation recommendation task involves a vastly larger and heavily
imbalanced candidate space. The high AUC values reflect the discriminability of
our features within the constrained evaluation setting, not the difficulty of the
full retrieval problem. The semantic ablation experiment (Section~\ref{sec:oa_sem})
directly addresses the concern that high performance may reflect topical redundancy.

\begin{table}[H]
\centering
\caption{5-Fold CV Results: OpenAlex Dataset (Mean $\pm$ Std)}
\label{tab:oa_cv}
\resizebox{\textwidth}{!}{
\begin{tabular}{lcccccc}
\toprule
\textbf{Model} & \textbf{ROC-AUC} & \textbf{PR-AUC} & \textbf{F1} & \textbf{Accuracy} & \textbf{MCC} & \textbf{Brier} \\
\midrule

Logistic Regression & $0.9677\pm0.0006$ & $0.9710\pm0.0006$ & $0.9073\pm0.0010$ & $0.9099\pm0.0010$ & $0.8211\pm0.0020$ & $0.0660\pm0.0006$ \\
Linear SVM & $0.9663\pm0.0005$ & $0.9698\pm0.0005$ & $0.9045\pm0.0009$ & $0.9068\pm0.0009$ & $0.8146\pm0.0017$ & $0.0680\pm0.0005$ \\
Random Forest & $0.9693\pm0.0006$ & $0.9702\pm0.0005$ & $0.9121\pm0.0011$ & $0.9119\pm0.0012$ & $0.8238\pm0.0024$ & $0.0646\pm0.0007$ \\
\textbf{Gradient Boosting} & $0.9762\pm0.0004$ & $0.9779\pm0.0004$ & $0.9219\pm0.0009$ & $0.9224\pm0.0009$ & $0.8449\pm0.0018$ & $0.0574\pm0.0006$ \\
\bottomrule
\end{tabular}}
\end{table}

\begin{figure}[H]
\centering
\includegraphics[width=0.85\textwidth]{../results/v2/figures/oa/oa_cv_v3.png}
\caption{Model performance comparison on the OpenAlex dataset (5-fold CV).}
\label{fig:oa_cv}
\end{figure}

\subsection{SHAP Feature Importance}

In the OpenAlex dataset, \texttt{prestige\_cited} is the
dominant predictor (mean $|$SHAP$|$ = 1.4615), followed by
\texttt{common\_citers} (1.2139).
The high SHAP value for prestige reflects the Matthew Effect operating strongly
in this corpus. Semantic similarity ranks second, confirming that topical relevance
is a major independent driver of citation behavior even when structural features
are available.

\begin{table}[H]
\centering
\caption{SHAP Feature Importance: OpenAlex Dataset (Gradient Boosting)}
\label{tab:oa_shap}
\begin{tabular}{lcc}
\toprule
\textbf{Feature} & \textbf{Mean $|$SHAP$|$} & \textbf{Rank} \\
\midrule

prestige\_cited & 1.4615 & 1 \\
semantic\_similarity & 1.2139 & 2 \\
common\_citers & 0.6833 & 3 \\
jaccard\_refs & 0.6370 & 4 \\
indegree\_velocity & 0.5593 & 5 \\
pagerank\_cited & 0.3726 & 6 \\
temporal\_gap & 0.3508 & 7 \\
common\_refs & 0.0861 & 8 \\
activity\_citing & 0.0140 & 9 \\
\bottomrule
\end{tabular}
\end{table}

\begin{figure}[H]
\centering
\includegraphics[width=0.75\textwidth]{../results/v2/figures/oa/oa_shap_v3.png}
\caption{SHAP feature importance for Gradient Boosting on the OpenAlex dataset.}
\label{fig:oa_shap}
\end{figure}

\subsection{Ablation Study and Semantic Ablation Experiment}
\label{sec:oa_sem}

Table~\ref{tab:oa_ablation} presents the leave-one-out ablation results for the
OpenAlex Gradient Boosting model. The most striking finding is the large AUC drop
when semantic similarity is removed: $-0.0268$,
by far the largest single-feature contribution. This finding directly addresses the
concern that the high overall AUC may reflect topical redundancy within the
constrained hard-negative setting: the model is not merely detecting topical
similarity, but semantic proximity is a genuinely dominant and independent predictor
of citation behavior.

\begin{table}[H]
\centering
\caption{Leave-One-Out Ablation: OpenAlex Dataset (Gradient Boosting)}
\label{tab:oa_ablation}
\begin{tabular}{lcc}
\toprule
\textbf{Feature Removed} & \textbf{AUC Without} & \textbf{AUC Drop} \\
\midrule

semantic\_similarity & 0.9495 & $-0.0268$ \\
prestige\_cited & 0.9734 & $-0.0029$ \\
common\_citers & 0.9734 & $-0.0029$ \\
temporal\_gap & 0.9740 & $-0.0023$ \\
pagerank\_cited & 0.9745 & $-0.0018$ \\
indegree\_velocity & 0.9761 & $-0.0002$ \\
activity\_citing & 0.9762 & $-0.0001$ \\
common\_refs & 0.9762 & $-0.0001$ \\
jaccard\_refs & 0.9762 & $-0.0001$ \\
\bottomrule
\end{tabular}
\end{table}

\begin{figure}[H]
\centering
\includegraphics[width=0.75\textwidth]{../results/v2/figures/oa/oa_ablation_v3.png}
\caption{Leave-one-out ablation study for Gradient Boosting on the OpenAlex dataset.}
\label{fig:oa_ablation}
\end{figure}

To further investigate the role of semantic similarity, we conducted a dedicated
semantic ablation experiment: all four models were trained and evaluated both with
and without the semantic similarity feature across all five folds. The results
confirm that removing semantic similarity reduces Gradient Boosting AUC from
0.9762 to 0.9494
(a drop of 0.0269). This demonstrates that structural
features alone achieve strong but not optimal performance, and that semantic
proximity provides a substantial, independent contribution.

\begin{figure}[H]
\centering
\includegraphics[width=0.75\textwidth]{../results/v2/figures/oa/oa_semantic_ablation_v3.png}
\caption{Semantic ablation experiment: model performance with and without semantic similarity (OpenAlex dataset).}
\label{fig:oa_sem_ablation}
\end{figure}

\subsection{Rolling Temporal Evaluation}

\begin{table}[H]
\centering
\caption{Rolling Temporal Evaluation: OpenAlex Dataset (ROC-AUC)}
\label{tab:oa_rolling}
\begin{tabular}{lcccc}
\toprule
\textbf{Test Window} & \textbf{LR} & \textbf{SVM} & \textbf{RF} & \textbf{GB} \\
\midrule

2005-2010 & 0.9609 & 0.9595 & 0.9569 & 0.9662 \\
2010-2015 & 0.9693 & 0.9681 & 0.9660 & 0.9744 \\
2015-2020 & 0.9732 & 0.9722 & 0.9705 & 0.9774 \\
2018-2025 & 0.9731 & 0.9722 & 0.9717 & 0.9788 \\
\bottomrule
\end{tabular}
\end{table}

\begin{figure}[H]
\centering
\includegraphics[width=0.85\textwidth]{../results/v2/figures/oa/oa_rolling_v3.png}
\caption{Rolling temporal evaluation for all models on the OpenAlex dataset.}
\label{fig:oa_rolling}
\end{figure}

Performance is stable and consistently high across all four test windows, with
Gradient Boosting achieving AUC above 0.966 in every period. The slight improvement
in later windows reflects the growing density of the citation network over time,
which enriches the structural features (prestige, PageRank, co-citation) and makes
the classification task progressively more tractable.

\subsection{Temporal Hold-Out Evaluation}

Training on pairs with citing year $\le 2015$ and testing on 2016--2020, Gradient
Boosting achieves AUC of 0.9799 on the OpenAlex
out-of-time test set, confirming robust temporal generalizability.


%% ============================================================
\section{Comparative Analysis: Dimensions vs.\ OpenAlex}
\label{sec:comp}
%% ============================================================

\subsection{Overall Predictive Performance}

Both datasets yield models with strong predictive performance, confirming that
citation dynamics in LIS are substantially predictable from historical network and
semantic states, regardless of the data source. Table~\ref{tab:crossdataset}
summarizes the key findings across the two datasets.

\begin{table}[H]
\centering
\caption{Summary of Cross-Dataset Findings}
\label{tab:crossdataset}
\begin{tabular}{lcc}
\toprule
\textbf{Characteristic} & \textbf{Dimensions} & \textbf{OpenAlex} \\
\midrule
Total pairs & 662,094 & 478,698 \\
Best model & Gradient Boosting & Gradient Boosting \\

Best CV AUC & 0.8949 & 0.9762 \\
Best Brier score & 0.1346 & 0.0871 \\
Hold-out AUC (2016--2020) & 0.8911 & 0.9799 \\
McNemar $p$ (RF vs GB) & $7.20e-26$ & $4.93e-14$ \\
Top SHAP feature & prestige\_cited & prestige\_cited \\
Semantic ablation drop & N/A & 0.0269 \\
\bottomrule
\end{tabular}
\end{table}

\begin{figure}[H]
\centering
\includegraphics[width=0.85\textwidth]{../results/v2/figures/comp/comparative_v3.png}
\caption{Comparative ROC-AUC across all models for both datasets.}
\label{fig:comp}
\end{figure}

The OpenAlex models consistently outperform their Dimensions counterparts. This
performance gap is attributable to two factors. First, the OpenAlex feature set
includes semantic similarity, which is the strongest single predictor in that
dataset. Second, the OpenAlex hard negatives are drawn from a tighter topical
neighborhood (primary topic rather than broad FoR code), which may make the task
slightly easier by reducing the diversity of negative examples.

\subsection{Feature Importance Hierarchy}

The most striking difference between the two datasets lies in the relative
importance of features. In the Dimensions dataset, prestige of the cited paper is
the dominant predictor, with co-citation overlap and in-degree velocity ranking
second and third. In the OpenAlex dataset, prestige remains the top structural
predictor, but semantic similarity is a close second with a much larger absolute
contribution.

Despite these differences, a consistent finding emerges: in both datasets, the
\textit{combination} of prestige, structural, and semantic features is necessary
for optimal performance. Neither prestige alone nor semantic similarity alone
achieves the performance of the full model, confirming the multi-dimensional nature
of citation behavior. The graph-theoretic features (PageRank, in-degree velocity)
contribute positively in both datasets, validating the enriched feature engineering
approach.

%% ============================================================
\section{Discussion}
\label{sec:discussion}
%% ============================================================

\subsection{The Predictability of Scientific Citations: Theoretical Implications}

Our results demonstrate that citation decisions in Library and Information Science
are substantially predictable from a small set of well-engineered features derived
from the historical state of the citation network. Gradient Boosting achieves
ROC-AUC of 0.895 (Dimensions) and 0.976 (OpenAlex) under strict temporal causality
constraints, with performance stable across four independent temporal test windows.
What does this high predictability mean for our understanding of scientific
knowledge production?

We argue that high citation predictability is not a trivial finding, nor is it
merely a methodological artifact of the balanced binary classification setting.
Rather, it reflects a fundamental property of knowledge diffusion within established
scientific disciplines: citations are not random acts of intellectual acknowledgment
but are \textit{structurally determined} by the position of papers in the existing
knowledge network. A paper is likely to be cited if it is highly visible (high
prestige and PageRank), if it is topically relevant to the citing paper (high
semantic similarity), if it is intellectually proximate (shared references,
co-citation overlap), and if it is temporally accessible (moderate age, recent
citation momentum).

This structural determinism connects to Kuhn's concept of normal science
\citep{kuhn1962structure}: within an established paradigm, the range of relevant
prior work is bounded by the shared assumptions and methods of the community.
Researchers working within a paradigm do not search the entire literature; they
search within the intellectual neighborhood defined by their research program.
This bounded search process is precisely what our features capture, and it is why
citations are so predictable. The high AUC values we observe are, in this reading,
a quantitative signature of paradigmatic closure.

This interpretation is reinforced by the dominance of prestige as a predictor.
Merton's Matthew Effect \citep{merton1968matthew} describes a positive feedback
mechanism in which recognition begets recognition. Our results confirm that this
mechanism operates strongly in LIS: papers that have already accumulated citations
are more likely to attract future ones, independent of their semantic content.
This is consistent with Crane's analysis of invisible colleges
\citep{crane1972invisible}: within a research community, the citation behavior of
individual researchers is shaped by the collective attention of the community,
which is itself concentrated on a small number of highly visible papers.

However, the strong independent contribution of semantic similarity (particularly
in the OpenAlex dataset, where removing it drops AUC by 0.027) suggests that the
Matthew Effect is not the only driver. Papers are not cited purely because they
are famous; they are cited because they are relevant. The interaction between
prestige and semantic proximity---captured by the non-linear ensemble models---
suggests that the most predictably cited papers are those that are both highly
visible and highly relevant: they occupy the center of the intellectual neighborhood
defined by the citing paper's research program.

The temporal features (temporal gap, in-degree velocity) add a further dimension:
citation probability is not static but evolves over a paper's life cycle. Papers
experiencing a recent surge in citations are more likely to continue attracting
citations, reflecting the dynamics of emerging research fronts where a new
contribution rapidly becomes a focal point for subsequent work \citep{redner2005citation}.
This temporal momentum is distinct from cumulative prestige and represents a genuine
additional signal about the current state of the knowledge diffusion process.

\subsection{The Sociology and Epistemology of Predictable Citations}

The high predictability of citations has deeper implications for the sociology and
epistemology of science. If citation behavior is substantially determined by
structural position, semantic proximity, and temporal momentum, then citation counts
are not pure measures of intellectual merit but are \textit{measures of structural
advantage}. A paper's citation count reflects not only its intrinsic quality but
also its position in the citation network at the time of its publication, the
prestige of its authors and venue, and the topical proximity of its content to the
mainstream of its field.

This has important implications for research evaluation. Citation-based metrics
such as the h-index or journal impact factor may systematically disadvantage
researchers working on novel or non-mainstream topics, or those outside of
well-established collaboration networks \citep{bornmann2008what}. Our finding that
co-citation overlap and bibliographic coupling are strong predictors of citation
suggests that papers embedded in dense citation neighborhoods---those that share
many references with existing literature---are more likely to be cited than papers
that introduce genuinely novel frameworks or methods that do not yet have established
citation neighborhoods.

Conversely, the strong role of semantic similarity suggests that citation metrics
do capture something genuine about intellectual relevance: papers are more likely
to be cited when they are topically relevant to the citing paper, independent of
their prestige. This finding supports a more nuanced view of citation metrics as
imperfect but not arbitrary measures of intellectual influence.

The practical implication is that citation prediction models like ours could be
used to construct ``expected citation'' baselines that account for the predictable
structural components of citation behavior, enabling fairer comparisons of research
impact across researchers with different career stages, network positions, and
topical niches. A paper that receives more citations than its structural position
predicts is, by this measure, genuinely influential; a paper that receives fewer
citations than predicted may be undervalued by the community.

\subsection{Network Science, Bibliometrics, and Machine Learning: A Synthesis}

Our study sits at the intersection of three traditions: network science, bibliometrics,
and machine learning. The enriched feature set, which includes year-capped PageRank
and in-degree velocity alongside traditional bibliometric features, demonstrates
that genuine graph-theoretic measures contribute incrementally to citation prediction
beyond raw citation counts. PageRank, which weights citations by the prestige of
the citing papers, provides a more nuanced measure of a paper's structural importance
than simple in-degree. In-degree velocity captures the temporal dynamics of citation
accumulation, distinguishing papers with steady long-term impact from those
experiencing a recent surge.

The modest but consistent contribution of journal prestige and author productivity
(Dimensions dataset) confirms that bibliometric metadata features add value beyond
the citation network structure itself. Journal prestige captures the reputation of
the publication venue, which influences the visibility of a paper to potential
citers. Author productivity captures the activity level of the citing team, which
may correlate with their breadth of reading and their tendency to cite a wider
range of papers.

The non-linear ensemble models (Random Forest, Gradient Boosting) consistently
outperform linear baselines (Logistic Regression, Linear SVM), confirming that
citation dynamics are non-linear and that the interactions between features are
important. The SHAP analysis reveals that these interactions are not merely
statistical artifacts but reflect genuine mechanisms: the benefit of high prestige
is amplified when semantic similarity is also high, and the penalty of large
temporal gap is mitigated when the cited paper has high in-degree velocity.

\subsection{Limitations}

Despite the robustness of our findings, our study has several limitations. First,
our analysis is restricted to Library and Information Science. The relative
importance of social, semantic, and prestige factors may vary substantially across
other disciplines, particularly the natural sciences, where citation practices
differ significantly. Second, while we substantially mitigate temporal leakage by
capping all features prior to the citing paper's publication year, observational
data cannot definitively establish causal mechanisms. Third, the OpenAlex dataset
has relatively low abstract coverage (64.3\%), which may introduce selection bias
in the semantic similarity feature. Fourth, our semantic similarity measure (TF-IDF
cosine) is less expressive than deep neural embeddings; future work should
investigate whether contextual embeddings provide additional predictive value while
maintaining leakage control. Fifth, the binary pair-classification evaluation
setting, while rigorous within its constraints, is a simplification of the real-world
citation recommendation task.

%% ============================================================
\section{Conclusion}
\label{sec:conclusion}
%% ============================================================

This paper demonstrates that machine learning approaches, when constrained by
strict temporal causality and enriched with a comprehensive feature set spanning
prestige, graph-theoretic network structure, bibliographic coupling, semantic
proximity, and bibliometric metadata, offer a powerful and interpretable framework
for understanding and predicting citation dynamics in Library and Information Science.

Our key empirical findings are: (i) Gradient Boosting achieves ROC-AUC of 0.895
(Dimensions) and 0.976 (OpenAlex) under 5-fold stratified cross-validation, with
well-calibrated probability estimates; (ii) performance is stable across four
independent temporal test windows, confirming temporal generalizability; (iii)
prestige of the cited paper is the dominant structural predictor in both datasets,
confirming the Matthew Effect; (iv) semantic similarity is the largest single-feature
contributor in OpenAlex (AUC drop of 0.027 when removed), confirming that topical
relevance is an independent driver of citation behavior; (v) graph-theoretic features
(PageRank, in-degree velocity) contribute meaningfully beyond raw citation counts.

Theoretically, we argue that high citation predictability reflects the bounded,
paradigmatic nature of knowledge diffusion within established disciplinary
communities: citations are structurally determined by a paper's position in the
existing knowledge network, its topical relevance, and its temporal momentum.
This structural determinism has important implications for research evaluation,
suggesting that citation-based metrics capture structural advantage as much as
intellectual merit, and motivating the development of expected-citation baselines
that account for the predictable components of citation behavior.

Future work should extend this framework to other disciplines, investigate the
temporal evolution of feature importance over a paper's life cycle, incorporate
richer graph-theoretic features (betweenness centrality, community membership),
and explore the use of contextual language model embeddings for semantic similarity
while maintaining rigorous leakage control.

\newpage
\bibliographystyle{plainnat}
\bibliography{references}

\end{document}
